% Generated draft: requires mathematical and source review.
\documentclass[11pt]{article}
\usepackage[margin=25mm]{geometry}
\usepackage{fontspec}
\setmainfont{Latin Modern Roman}
\usepackage{amsmath,amssymb,mathtools,unicode-math}
\setmathfont{Latin Modern Math}
\usepackage{enumitem,needspace,xurl,hyperref}
\setlength{\parindent}{0pt}
\setlength{\parskip}{5pt}
\setlength{\emergencystretch}{3em}
\newcommand{\N}{\mathbb{N}}
\newcommand{\Z}{\mathbb{Z}}
\newcommand{\Q}{\mathbb{Q}}
\newcommand{\R}{\mathbb{R}}
\newcommand{\C}{\mathbb{C}}
\newcommand{\F}{\mathbb{F}}
\newcommand{\A}{\mathbb{A}}
\newcommand{\PP}{\mathbb{P}}
\newcommand{\E}{\mathbb{E}}
\newcommand{\eps}{\varepsilon}
\newcommand{\dd}{\,\mathrm{d}}
\begin{document}
\section*{UnsolvedMath ID: 30004499}
\textbf{Source problem code:} catalogue entry
\subsection{Definitions and mathematical statement}
Let \(\A_{\Q}\) denote the adele ring of \(\Q\), and let \(Z(G)\) denote the center of a group \(G\). Using the symbols in the supplied record, put
\[
G=D^{\times}(\A_{\Q}),\qquad
X=D^{\times}(\Q)\backslash G/\bigl((K_\infty\cap G)Z(G)\bigr).
\]
The record does not specify \(D\), the ambient group containing \(K_\infty\), or the action relevant to the representation condition below. These symbols therefore retain their unresolved source meanings.

The supplied hypothesis is that there are no nonzero functions on \(X\) transforming according to the stated \(3\)-dimensional spherical representation of \(\operatorname{SO}(3)\). Denote this hypothesis by \(H\), and denote the unspecified collection of special eigenfunctions by \(\mathcal E\). The cleaned statement asks whether
\[
H\ \Longrightarrow\
\exists\,\mathcal E_0\subseteq\mathcal E,\quad
|\mathcal E_0|<\infty,\quad
\forall f\in\mathcal E\setminus\mathcal E_0\quad
\forall x\in G,\quad
\text{\(x\) is a local maximum point of \(f\)}.
\]
This is only a transcription of the cleaned formulation: the original statement instead quantifies over ``every such point \(x\)'', whose antecedent is absent from the supplied continuation fragment. Consequently, the intended set of points may be a previously defined subset rather than all of \(G\).

For a real-valued function \(f\) on a topological space \(M\), a point \(x\in M\) is a local maximum point if there exists a neighborhood \(U\) of \(x\) such that \(f(y)\leq f(x)\) for every \(y\in U\). The record does not identify the domain or codomain of the special eigenfunctions, the topology in which maxima are taken, or whether representatives \(x\in G\) are to be evaluated through a quotient map. If the eigenfunctions are complex-valued, it also does not specify an ordered real-valued quantity, such as their absolute value, to maximize. No such interpretation is imposed here. The missing context prevents a fully specified mathematical statement.
\subsection{Short English statement}
Under a stated vanishing condition for functions transforming according to a \(3\)-dimensional spherical representation of \(\operatorname{SO}(3)\), the question asks whether every relevant point is a local maximum for all but finitely many special eigenfunctions. The eigenfunctions, representation action, and relevant points require missing context; the original wording does not justify replacing ``every such point'' by every adelic group element.
\subsection{Sources}
\begin{itemize}
\item Supplied local UnsolvedMath registry record, permanent ID \(30004499\), title \emph{Local Maxima of Special Eigenfunctions at Spherical Points}, source code \texttt{catalogue entry}, in the Oberwolfach Reports Open Problems collection. The record gives proposed year \(2020\) and names no proposer.
\item Supplied source citation: \emph{Automorphic Forms and Arithmetic (2021)}, Oberwolfach Reports, pp.\ 1317--1383, DOI \texttt{10.4171/owr/2020/26}. Source URL and supplied literature reference, labeled \emph{Oberwolfach Report 26/2020}: \url{https://doi.org/10.4171/owr/2020/26}. The citation's year \(2021\), report designation \(26/2020\), and proposed year \(2020\) are retained as supplied.
\item The registry labels the statement \texttt{corrected\_verified} and reports that the official report was checked against Question 3 to restore the adelic quotient, spherical representation, finite exceptional set, and local-maximum question. This is supplied verification metadata, not an independent verification performed for this draft.
\item Literature assessment dated exactly \texttt{2026-08-22}: supplied dated curation labels the problem open and describes the extracted text as a continuation fragment concerning conditional methods for period integrals and a local-maximum question, with no proof of the underlying general claim recorded. The background identifies Oberwolfach Reports as workshop reports containing open-problem discussions. The supplied assessment expressly requires independent verification of time-sensitive status and cited claims; it is not independently verified here.
\end{itemize}
\end{document}
